% ══════════════════════════════════════════════
% CHAPTER 2: ARCHITECTURE & PROJECT STRUCTURE
% ══════════════════════════════════════════════

\chapter{Architecture \& Project Structure}

\section{Technology Stack}

The Active Dream Bot is built with a carefully selected, modern technology stack optimized for reliability, performance, and maintainability.

\begin{table}[H]
\centering
\caption{Technology Stack}
\begin{tabularx}{\textwidth}{|l|l|X|}
\hline
\textbf{Component} & \textbf{Technology} & \textbf{Why This Choice} \\
\hline
Language & Python 3.11+ & Industry standard for bot development; async support; huge ecosystem \\
\hline
Bot Framework & python-telegram-bot v21.6 & Official, well-maintained; fully async; 40K+ GitHub stars \\
\hline
Database & SQLite (via aiosqlite) & Zero configuration; serverless; file-based; perfect for single-bot deployments \\
\hline
SMS API & 5sim.net & Reliable; 200+ countries; fair pricing; great API documentation \\
\hline
HTTP Client & aiohttp 3.10 & High-performance async HTTP; used for 5sim API calls \\
\hline
Configuration & python-dotenv & Secure environment variable management via .env files \\
\hline
\end{tabularx}
\end{table}

\begin{infobox}
All libraries are open-source and free. The only paid service is the 5sim.net API for purchasing virtual numbers.
\end{infobox}

% ──────────────────────────────────────────────
\section{Complete Project Structure}

The project follows a clean, modular architecture. Every component has a single responsibility.

\begin{lstlisting}[style=bashstyle, title=Project Directory Tree]
active_dream_bot/
|
|-- bot.py                  # Main entry point - starts the bot
|-- config.py               # Loads configuration from .env
|-- .env                    # YOUR secrets (never share this!)
|-- .env.example            # Template for .env
|-- requirements.txt        # Python dependencies
|-- README.md               # Quick-start guide
|-- active_dream.db         # SQLite database (auto-created)
|
|-- database/
|   |-- __init__.py
|   |-- db.py               # Database connection & table creation
|   |-- models.py           # All database query functions
|
|-- handlers/
|   |-- __init__.py
|   |-- start.py            # /start command & access request flow
|   |-- menu.py             # Main menu & back navigation
|   |-- numbers.py          # Get number, OTP polling, delete number
|   |-- balance.py          # Balance display & withdrawal
|   |-- admin.py            # Admin panel, approve/reject, broadcast
|
|-- services/
|   |-- __init__.py
|   |-- sms_provider.py     # Abstract interface for SMS APIs
|   |-- fivesim.py          # 5sim.net API implementation
|
|-- utils/
|   |-- __init__.py
|   |-- constants.py        # Messages, service/country definitions
|   |-- keyboards.py        # All inline keyboard layouts
|   |-- decorators.py       # Permission check decorators
|
|-- docs/                   # This documentation
\end{lstlisting}

% ──────────────────────────────────────────────
\section{File-by-File Explanation}

\subsection{\texttt{bot.py} --- Main Entry Point (105 lines)}

This is the file you run to start the bot. It:
\begin{itemize}
    \item Validates configuration (checks if all required .env values are set)
    \item Initializes the database (creates tables if they don't exist)
    \item Registers all command handlers (\texttt{/start}, \texttt{/admin}, \texttt{/addbal}, etc.)
    \item Registers all callback query handlers (button presses)
    \item Starts the Telegram polling loop
    \item Notifies admins when the bot goes online
    \item Cleanly shuts down the database on exit
\end{itemize}

\subsection{\texttt{config.py} --- Configuration Loader (42 lines)}

Reads all settings from the \texttt{.env} file using \texttt{python-dotenv}. Includes a \texttt{validate()} method that checks for missing required values before the bot starts. This prevents runtime errors from missing API keys.

\subsection{\texttt{database/db.py} --- Database Engine (72 lines)}

Manages the SQLite database connection using \texttt{aiosqlite} (async SQLite). Creates all tables on first run:
\begin{itemize}
    \item \textbf{users} --- Stores all user records with approval status
    \item \textbf{numbers} --- Stores all number purchase records
    \item \textbf{transactions} --- Stores all balance changes (deposits, purchases, refunds)
    \item \textbf{settings} --- Key-value store for bot configuration
\end{itemize}

Uses WAL (Write-Ahead Logging) journal mode for better concurrent read performance.

\subsection{\texttt{database/models.py} --- Database Queries (210 lines)}

Contains all database operations organized into three sections:
\begin{enumerate}
    \item \textbf{User Operations:} Create user, get user, update status, get all users, count users
    \item \textbf{Balance Operations:} Get balance, add balance, deduct balance, get transactions
    \item \textbf{Number Operations:} Create number record, update status, get active numbers, get history
\end{enumerate}

\subsection{\texttt{handlers/start.py} --- Access Control (130 lines)}

Implements the complete permission flow:
\begin{itemize}
    \item \texttt{start\_command()} --- Checks user status and shows appropriate screen
    \item \texttt{request\_access\_callback()} --- Handles the ``Request Access'' button
    \item \texttt{check\_status\_callback()} --- Lets pending users check their status
\end{itemize}

\subsection{\texttt{handlers/numbers.py} --- Number \& OTP Flow (340 lines)}

The largest and most complex handler. Manages:
\begin{itemize}
    \item Service selection (Facebook, Instagram, etc.)
    \item Country selection (20+ countries)
    \item Number purchase via 5sim.net API
    \item \textbf{Background OTP polling} --- Uses \texttt{asyncio.create\_task()} to poll every 5 seconds
    \item OTP display and group forwarding
    \item Number deletion and balance refunds
    \item Timeout handling (auto-cancel after 5 minutes)
\end{itemize}

\subsection{\texttt{handlers/admin.py} --- Admin Panel (260 lines)}

Full admin management system:
\begin{itemize}
    \item Dashboard with live statistics
    \item Pending request list
    \item User approval/rejection
    \item Balance management
    \item Broadcasting to all users
    \item Ban/Unban functionality
\end{itemize}

\subsection{\texttt{services/sms\_provider.py} --- Abstract Interface (50 lines)}

Defines the contract that any SMS API provider must implement. This makes the bot \textbf{provider-agnostic} --- you can switch from 5sim.net to sms-activate.org or any other service by implementing this interface.

Methods required:
\begin{itemize}
    \item \texttt{get\_balance()} --- Check API account balance
    \item \texttt{buy\_number(country, service)} --- Purchase a virtual number
    \item \texttt{check\_sms(order\_id)} --- Check if OTP has arrived
    \item \texttt{cancel\_number(order\_id)} --- Cancel/release a number
    \item \texttt{finish\_number(order\_id)} --- Mark order as complete
    \item \texttt{get\_prices(country, service)} --- Get pricing info
\end{itemize}

\subsection{\texttt{services/fivesim.py} --- 5sim.net Implementation (170 lines)}

Implements the \texttt{SmsProvider} interface for 5sim.net. Handles:
\begin{itemize}
    \item JWT authentication with the 5sim API
    \item Number purchasing with automatic operator selection
    \item SMS polling with status parsing (PENDING, RECEIVED, CANCELED, TIMEOUT)
    \item OTP code extraction from SMS messages
    \item Error handling for all API calls
\end{itemize}

\subsection{\texttt{utils/keyboards.py} --- Keyboard Layouts (170 lines)}

All Telegram inline keyboard definitions in one place:
\begin{itemize}
    \item Request access keyboard
    \item Main menu keyboard (4 buttons + channel link)
    \item Service selection keyboard (13 services, 2 per row)
    \item Country selection keyboard (20 countries, 2 per row)
    \item Number active/completed keyboards
    \item Admin panel keyboard
    \item Admin user action keyboard
\end{itemize}

\subsection{\texttt{utils/decorators.py} --- Permission Decorators (40 lines)}

Two Python decorators that wrap handler functions:
\begin{itemize}
    \item \texttt{@require\_approved} --- Only allows approved users to access the handler
    \item \texttt{@require\_admin} --- Only allows admin users (defined in .env)
\end{itemize}

These are applied to every handler that needs protection, ensuring no unauthorized access.

% ──────────────────────────────────────────────
\section{Database Schema}

\subsection{Users Table}
\begin{table}[H]
\centering
\begin{tabularx}{\textwidth}{|l|l|l|X|}
\hline
\textbf{Column} & \textbf{Type} & \textbf{Default} & \textbf{Description} \\
\hline
user\_id & INTEGER (PK) & --- & Telegram user ID \\
username & TEXT & NULL & Telegram @username \\
first\_name & TEXT & NULL & User's first name \\
last\_name & TEXT & NULL & User's last name \\
status & TEXT & `new' & new / pending / approved / rejected / banned \\
balance & REAL & 0.0 & User's credit balance in USD \\
created\_at & TIMESTAMP & now & First interaction time \\
approved\_at & TIMESTAMP & NULL & When admin approved \\
approved\_by & INTEGER & NULL & Admin's user ID \\
\hline
\end{tabularx}
\end{table}

\subsection{Numbers Table}
\begin{table}[H]
\centering
\begin{tabularx}{\textwidth}{|l|l|l|X|}
\hline
\textbf{Column} & \textbf{Type} & \textbf{Default} & \textbf{Description} \\
\hline
id & INTEGER (PK) & auto & Auto-increment ID \\
user\_id & INTEGER (FK) & --- & Who requested this number \\
phone\_number & TEXT & --- & The virtual phone number \\
country & TEXT & --- & Country name \\
country\_code & TEXT & --- & Flag emoji \\
service & TEXT & --- & Service name (Facebook, etc.) \\
order\_id & TEXT & --- & 5sim order ID \\
status & TEXT & `active' & active / waiting / completed / cancelled / expired \\
sms\_code & TEXT & NULL & Extracted OTP code \\
sms\_full & TEXT & NULL & Full SMS text \\
cost & REAL & 0.0 & Cost of this number \\
created\_at & TIMESTAMP & now & When number was assigned \\
completed\_at & TIMESTAMP & NULL & When OTP was received/cancelled \\
\hline
\end{tabularx}
\end{table}

\subsection{Transactions Table}
\begin{table}[H]
\centering
\begin{tabularx}{\textwidth}{|l|l|l|X|}
\hline
\textbf{Column} & \textbf{Type} & \textbf{Default} & \textbf{Description} \\
\hline
id & INTEGER (PK) & auto & Auto-increment ID \\
user\_id & INTEGER (FK) & --- & User ID \\
amount & REAL & --- & Amount (positive = credit, negative = debit) \\
type & TEXT & --- & deposit / purchase / refund / admin\_add / withdrawal \\
description & TEXT & --- & Human-readable description \\
created\_at & TIMESTAMP & now & Transaction time \\
\hline
\end{tabularx}
\end{table}

% ──────────────────────────────────────────────
\section{System Architecture Diagram}

\begin{center}
\begin{tikzpicture}[
    node distance=1.5cm,
    box/.style={rectangle, draw=primaryblue, fill=primaryblue!10, rounded corners, minimum width=3cm, minimum height=1cm, text centered, font=\small\bfseries},
    db/.style={cylinder, draw=successgreen, fill=successgreen!10, shape border rotate=90, aspect=0.3, minimum width=2cm, minimum height=1.2cm, text centered, font=\small\bfseries},
    api/.style={rectangle, draw=warningyellow, fill=warningyellow!10, rounded corners, minimum width=3cm, minimum height=1cm, text centered, font=\small\bfseries},
    arrow/.style={->, thick, >=stealth}
]
    % Nodes
    \node[box] (user) {Telegram User};
    \node[box, right of=user, xshift=3cm] (bot) {Active Dream Bot};
    \node[db, below of=bot, yshift=-1cm] (db) {SQLite DB};
    \node[api, right of=bot, xshift=3cm] (fivesim) {5sim.net API};
    \node[box, below of=user, yshift=-1cm] (admin) {Admin User};
    \node[box, below of=fivesim, yshift=-1cm] (group) {OTP Group};
    
    % Arrows
    \draw[arrow] (user) -- node[above, font=\tiny] {/start, buttons} (bot);
    \draw[arrow] (bot) -- node[right, font=\tiny] {read/write} (db);
    \draw[arrow] (bot) -- node[above, font=\tiny] {buy, check SMS} (fivesim);
    \draw[arrow] (admin) -- node[above, font=\tiny] {approve/reject} (bot);
    \draw[arrow] (bot) -- node[above, font=\tiny] {forward OTP} (group);
    \draw[arrow, dashed] (fivesim) -- node[right, font=\tiny] {number + OTP} (bot);
\end{tikzpicture}
\end{center}

\begin{infobox}
The architecture is intentionally simple --- a single Python process handles everything. For high-traffic scenarios (100+ concurrent users), consider upgrading to PostgreSQL and deploying with webhook mode instead of polling.
\end{infobox}
